\chapterhead{1}{A Cup That Seems to Lose Water}{Marking a level and finding out where the water went}\label{ch1}

A puddle on a sunny path is gone by afternoon. A cup of water on a shelf has a little less water each day. Where did the water go?

There are three ideas. The water might have \textbf{soaked away} into something. It might have \textbf{stopped existing}. Or it might have \textbf{moved} somewhere we cannot easily see. A good test does not just pick the idea we like. It sets one idea against another and sees what happens.

\begin{center}
\begin{tikzpicture}
  \foreach \x/\lab in {0/open cup,5/covered cup}{
    \draw[line width=1.4pt,color=ink] (\x,0) -- (\x,2.2) (\x+1.5,0) -- (\x+1.5,2.2) (\x,0) -- (\x+1.5,0);
    \fill[sky!50] (\x+0.03,0.03) rectangle (\x+1.47,1.4);
    \draw[color=tomato,line width=1.2pt] (\x-0.2,1.4) -- (\x+1.7,1.4);
    \node[font=\sffamily\scriptsize] at (\x+0.75,-0.35) {\lab};
  }
  \draw[line width=1.6pt,color=ink!70] (4.9,2.28) -- (6.6,2.28);
  \node[font=\sffamily\scriptsize,anchor=west] at (7.0,1.4) {the red line marks the starting level};
\end{tikzpicture}
\end{center}

\begin{tryit}{Mark the level}
\textbf{You need:} two clear cups that are the same, water, a marker pen, a ruler, cling film, tape.
\begin{enumerate}[leftmargin=1.6em]
\item Pour the same amount of water into both cups, about 6~cm deep. Use the ruler and pour until each is exactly the same.
\item Draw a line on the outside at the water level of each cup.
\item Cover \textbf{one} cup tightly with cling film and tape. Leave the other open.
\item Stand both on the same shelf, side by side, out of direct sun. Measure the water depth in millimetres at the same time each day for five days.
\end{enumerate}
\par\medskip
\renewcommand{\arraystretch}{1.6}
\begin{tabular}{|l|c|c|c|c|c|c|}
\hline
\textbf{Day} & 0 & 1 & 2 & 3 & 4 & 5\\ \hline
\textbf{Open cup (mm)} & & & & & & \\ \hline
\textbf{Covered cup (mm)} & & & & & & \\ \hline
\end{tabular}
\end{tryit}
\adult{Use cold tap water. Put the cups where they will not be knocked over. Do not drink the water afterwards.}

\subhead{Reading a level honestly}
Water curves a little where it touches the glass. Read the flat middle of the surface, with your eye level with it, and round to the nearest millimetre. A one-millimetre wobble in a reading is not a change in the water. We saw why in Book~1.

\begin{worked}{Practice numbers, not from a real test}
Open cup, in mm: 60, 58, 56, 53, 51, 49. Covered cup: 60, 60, 60, 59, 60, 60.\\
The open cup dropped $60-49=11$~mm in 5 days. That is $11\div5=2.2$~mm a day, or about 2~mm.\\
The covered cup read 59 once and 60 at the next reading, so the 59 is a wobble, not a loss.
\end{worked}

\begin{center}
\begin{tikzpicture}
  \node[bigcir,minimum size=2.2cm] (a) at (0,0) {}; \node[font=\sffamily\scriptsize,align=center] at (0,0) {open cup\\60 mm};
  \node[valcir] (b) at (4.6,0) {49 mm};
  \draw[popA] (a) -- node[above,note] {5 days} (b);
  \node[wrapcir,minimum size=2.6cm] (w) at (8.2,0) {}; \node[valcir] at (8.2,0.35) {$-\,11$}; \node[note] at (8.2,-0.4) {mm left the cup};
  \node[note,anchor=west,text width=3cm] at (9.7,0) {\textbf{kept:} the drop\\\textbf{set aside:} the path the water took};
\end{tikzpicture}
\end{center}

\begin{multix}[hmm]\multixbot[hmm]\tcblower
The open cup dropped 11 mm and the covered cup did not, so the extra water just stopped existing. I will write that in my notebook. Should I?
\end{multix}

\trust{We changed one thing, the cling film, and kept the cups, the water, the shelf and the start level the same, then measured both cups every day.}
{The covered cup held its level and the open one dropped steadily, so the lid made the difference.}
{Where the missing water went, or whether the same drop would happen in a cold room or in the sun. We tested one shelf.}

\begin{keepbox}
We turned a cup of water into a number of millimetres. \textbf{We kept} \wline[3.4cm] \textbf{and we set aside} \wline[3.4cm].
\end{keepbox}

\clearpage

\begin{puzzle}{Which cup lost water?}
Open cup: 40, 38, 37, 35~mm. Covered cup: 40, 40, 40, 40~mm.\\
How much did the open cup drop? \wline[1.6cm]~mm \quad The covered cup? \wline[1.6cm]~mm
\end{puzzle}

\begin{puzzle}{Water a day}
An open cup read 50, 47, 45, 42, 40~mm over four days after the first reading.\\
Total drop: \wline[1.6cm]~mm \quad Average per day: \wline[1.6cm]~mm
\end{puzzle}

\begin{puzzle}{Predict the level}
At 2~mm a day, a cup that starts at 50~mm will read what after 10 days? \wline[1.8cm]~mm\\
Why is that only a guess? \wline[5cm]
\end{puzzle}

\begin{puzzle}{Wobble or loss?}
A covered cup read 60, 60, 59, 60, 60~mm. Did it lose water? \wline[2cm]\\
What makes you think so? \wline[5cm]
\end{puzzle}

\begin{puzzle}{Spot the unfair test}
Ana put her open cup on a sunny window sill and her covered cup in a dark cupboard. The open cup lost water. She says the lid saved it. What is wrong? \wline[6cm]
\end{puzzle}

\clearpage

\begin{puzzle}{Three ideas}
The water might have \emph{soaked away}, \emph{stopped existing} or \emph{moved}. The covered cup kept all its water. Which idea does that make less likely for the open cup, and why? \wline[6cm]
\end{puzzle}

\begin{puzzle}{Draw the move}
The open cup started at 60~mm and ended at 49~mm. Fill in the empty circles.
\begin{center}
\begin{tikzpicture}
  \node[cir] (a) at (0,0) {60}; \node[cir,dashed] (b) at (5,0) {}; 
  \draw[moveA] (a) -- node[above,note] {11 mm} (b);
  \node[note] at (0,-0.8) {cup before}; \node[note] at (5,-0.8) {air gained?};
  \node[cir,dashed] (c) at (9,0) {}; \node[note] at (9,-0.8) {cup after};
\end{tikzpicture}
\end{center}
\end{puzzle}

\begin{puzzle}{A puddle}
One puddle in the sun is gone by noon. A puddle in the shade of a wall is still there at night. What is different about the two puddles? \wline[5cm]
\end{puzzle}

\begin{puzzle}{Around you}
Name two places at home where water seems to go away by itself.\\
\wline[5cm]\\[8pt]\wline[5cm]
\end{puzzle}

\begin{puzzle}{Missing water}
A student says ``the water in the open cup went through the glass''. If that were true, what would happen to the covered cup? \wline[6cm]
\end{puzzle}

\clearpage
\thispagestyle{fancy}
\subhead{Cup diary}
Keep your own record. Fill in what you see each day, and draw the level on the cups.\par\medskip
\renewcommand{\arraystretch}{3.0}
\begin{tabularx}{\linewidth}{|>{\sffamily\small\bfseries}p{1.4cm}|X|X|X|}
\hline
Day & Open cup (mm) & Covered cup (mm) & What I saw on the cling film \\ \hline
0 & & & \\ \hline
1 & & & \\ \hline
2 & & & \\ \hline
3 & & & \\ \hline
4 & & & \\ \hline
5 & & & \\ \hline
\end{tabularx}

\vspace{6pt}
Do you see anything on the inside of the cling film? Draw it and write about it. You will see why it matters in Chapter~3.\par\medskip
\begin{tikzpicture}\draw[ink!70,line width=1pt,rounded corners=3pt] (0,0) rectangle (11.4,5);\end{tikzpicture}
\clearpage
